% ----------------------------------------------------------
% Metodologia
% ----------------------------------------------------------

\chapter{Trabalhos Relacionados}
No seguinte capítulo serão discutidos os pontos relacionados a sistemas semelhantes ao proposto, bem como suas funcionalidades e características. 

\section{CALCULADORA DE PROGRESSÃO DA PENA}
No presente sistema web, o seu principal objetivo é trazer uma estimativa de quando um condenado poderá mudar para um regime menos rigoroso, baseado na Lei de Execução Penal  e nas respectivas alterações do Pacote Anticrime:\cite{jurisdiques}

\begin{itemize}
    \item Cálculo automático da progressão de regime;
    \item Base de cálculo de acordo com o LEP e pacote anti-crime;
    \item Parâmetros editáveis(Pena total e pena já cumprida, tipo de crime, condição do condenado); e
    \item Percentuais de cumprimento de pena;
\end{itemize}

Contudo, o sistema desenvolvido possui algumas limitações: não faz menção à integração com cálculos de remição de pena por trabalho, estudo e leitura; não considera a unificação de penas; depende do usuário para a inserção de dados corretos; não avalia requisitos subjetivos.

No sistema desenvolvido em questão, além das operações básicas para base de cálculo, terão a inserção de parâmetros essenciais para a execução do procedimento, como: campo para inserção de remição de pena, inserção de unificação de pena; os dados são inseridos manualmente pelo usuário, com previsão futura 
de extração automática a partir de documentos oficiais.

\section{CALCULADORA DE PRESCRIÇÃO DA PRETENSÃO EXECUTÓRIA}
O sistema web  calculadora de prescrição da pretensão executória, disponibilizado pelo Conselho Nacional de Justiça (CNJ),  tem o intuito de realizar o  controle de prazos prescricionais e tem-se as seguintes funcionalidades:\cite{cnjCalculadora}

\begin{itemize}
\item Permite calcular prazos da prescrição da pretensão punitiva (que ocorre antes do trânsito em julgado) e executória (ocorre após condenação definitiva);

\item Possui campos para informações do acusado, número do processo, dados referente a data da denúncia, sentença;

\item Considera causas de aumentos ou diminuições de penas; e

\item Dados referentes a data provável da prescrição; 
\end{itemize}

Entretanto, o sistema desenvolvido, apesar de possuir muitas ferramentas para fins de cálculo de prescrição e ser uma ferramenta desenvolvida principalmente para advogados, magistrados e servidores, apresenta uma interface pouco intuitiva, não realiza integração com sistemas processuais, não gera relatório com a finalidade de comprovar as atividades avaliadas do acusado.

Pensando nesses pontos, por mais que no sistema apresentado o seu principal foco seja a prescrição penal, ele possui características que se assemelham ao projeto proposto, que facilita o dia a dia de usuários que têm a necessidade de ter uma ferramenta adequada para cumprimento do processo burocrático de execução penal. Dentre suas funcionalidades básicas, tem como foco o cálculo de execução penal, com os elementos que podem diminuir (remição) ou aumentar (regressão) o tempo de pena de um acusado, tendo como a base principal de fundamentação a LEP.

\section{CALCULADORA DE EXECUÇÃO PENAL}
O seguinte sistema web Calculadora de Execução Penal, desenvolvido pela Equipe Busca CDP, tem como objetivo principal realizar o cálculo de execução penal, com foco na progressão de regime, baseado na lei do pacote anti crime. Dessa forma, possui as seguintes funcionalidades:\cite{buscacdp}

\begin{itemize}
    \item Calcula tempo total da pena;
    \item Calcula a progressão de regime;
    \item Realiza o cálculo de livramento condicional;
    \item Possui os campos para parâmetros editáveis;
    \item Realiza o cálculo para remição por trabalho ou estudo
    \item Baseado na lei do pacote anti crime 13.964/2019 \cite{pacoteAnticrime}
\end{itemize}

No entanto, a ferramenta, apesar de ser intuitiva e abranger os benefícios de progressão, livramento condicional e remição, possui algumas limitações: não considera a unificação de penas e detração, não gera um documento acerca das informações calculadas, não realiza o armazenamento de informações para consultas futuras.

Nessa perspectiva, mesmo que o foco do sistema exposto seja realizar o cálculo completo de execução penal, ele ainda apresenta ausência de parâmetros que são necessários para a construção do cálculo completo para uma execução penal. Dessa forma, o sistema desenvolvido neste projeto abrange todos os parâmetros necessários para que se realize o cálculo integral de execução penal, incluindo parâmetros como detração, remição por leitura, unificação de penas e o armazenamento dos dados, baseando-se na lei de execução penal com as alterações feitas pela lei do pacote anti crime.\cite{lep, pacoteAnticrime}

\section{CALCULADORA DE EXECUÇÃO PENAL}
O presente trabalho tem como principal objetivo realizar conversões de dias para anos e vice-versa, bem como cálculos básicos de soma e subtração de datas, de acordo com as frações de lapso temporal da LEP. Logo, estas são as funcionalidades presentes no sistema desenvolvido:\cite{fernandoCruzCalc, lep}

\begin{itemize}
    \item Conversão de tempo;
    \item Frações; e
    \item Porcentagens.
\end{itemize}

Entretanto, para fins de cálculos de execução penal, em que se necessita dos parâmetros de progressão, remição, detração e unificação de pena, o sistema exposto apresenta limitações quanto ao objetivo deste projeto. O sistema pode ser usado como fonte à parte para o desenvolvimento de um processo de execução, visto que ele realiza as conversões necessárias, com a finalidade de otimizar tempo.


\section{CALCULADORA DE EXECUÇÃO PENAL ONLINE}
A calculadora de Execução Penal Online segue os mesmos princípios da calculadora de execução penal, apresentada anteriormente. Ela configura-se para realizar conversões de tempo e frações, que também age de acordo com a Lei de Execução Penal.\cite{saraBarbozaCalc}

\begin{itemize}
    \item Conversão de tempo; 
    \item Interface amigável; e
    \item Frações.
\end{itemize}

Diante do exposto, observa-se que a ferramenta de Execução Penal online apresenta uma limitação quanto às suas finalidades. Apesar de oferecer poucas ferramentas, a interface do sistema é intuitiva e agradável. Dessa forma, o CAlPEC facilita e otimiza o tempo do usuário, fazendo com que ele encontre as ferramentas necessárias em um único local.

\section{TRABALHOS RELACIONADOS}

O Sistema CalPEC foi desenvolvido com o principal intuito de facilitar e otimizar o tempo que é levado em consideração para a realização de cálculos contínuos de execução penal, pensando nos usuários que precisam administrar diariamente mais de um processo ou até mesmo para aqueles que usam a ferramenta como objeto de estudo. Enquanto os sistemas já desenvolvidos ficam em torno de cálculos simples (conversões) ou ficam restritos a poucos parâmetros (progressão e/ou remição), a calculadora de execução criminal une as principais funcionalidades em um único sistema.

Diferente da Calculadora de progressão da pena, cujo principal objetivo é trazer uma estimativa de quando um condenado poderá mudar de regime, com percentual atual e tempo restante para cumprir o tempo necessário para progressão; ou da Calculadora de prescrição da pretensão executória, que tem o intuito de realizar o controle de prazos prescricionais, a CalPEC destaca-se por oferecer os parâmetros necessários para progredir de regime. \cite{cnjCalculadora, buscacdp}

Do mesmo modo, a Calculadora de execução penal e a Calculadora de execução penal online assemelham-se por terem praticamente o mesmo intuito: garantir a conversão em dias, meses ou anos de uma determinada data. Apesar de suas funcionalidades principais serem atrativas e necessárias, pois é usada constantemente a conversão de datas, essas ferramentas ficam limitadas a oferecer somente isso, o que provoca no usuário uma insatisfação do produto, já que, caso ele queira optar por calcular uma remição, não será possível, pois os sistemas não apresentam essa opção. Ao contraposto, a CalPEC possui as opções de calcular remição ou progressão, estando disponíveis para o usuário sempre que ele precisar. \cite{fernandoCruzCalc, saraBarbozaCalc}

A ferramenta Calculadora de progressão da pena assemelha-se a CalPEC, pois ambas oferecem a progressão de regime. No entanto, a Calculadora de progressão da pena apresenta como foco a progressão, disponibilizando o percentual exigido e tempo necessário para progredir de regime. 

Em resumo, o CalPEC destacou-se em relação às demais ferramentas, pois ofereceu solução aos problemas identificados, reunindo as principais funcionalidades em um único sistema e tornando a experiência do usuário um diferencial. A Tabela~\ref{tab:calculadoras} sintetiza a comparação entre as ferramentas analisadas e o CalPEC, evidenciando as funcionalidades presentes em cada uma.

\clearpage
\begin{table}[h]
\centering
\caption{Trabalhos relacionados e suas funcionalidades}
\label{tab:calculadoras}
\resizebox{\textwidth}{!}{
\begin{tabular}{|l|c|c|c|c|c|}
\hline
\textbf{Funcionalidade} & \textbf{Cal. Progressão} & \textbf{Calc. Prescrição CNJ} & \textbf{Calc. Execução Penal} & \textbf{Calc. de Pena} & \textbf{CalPEC} \\
\hline
Cadastro de Apenados & \ding{55} & \ding{51} & \ding{55} & \ding{55} & \ding{51} \\
\hline
Progressão de Regime & \ding{51} & \ding{55} & \ding{51} & \ding{51} & \ding{51} \\
\hline
Remição de Pena & \ding{55} & \ding{51} & \ding{51} & \ding{55} & \ding{51} \\
\hline
Detração de Pena & \ding{55} & \ding{51} & \ding{55} & \ding{55} & \ding{51} \\
\hline
Alertas de Progressão Vencida & \ding{55} & \ding{55} & \ding{55} & \ding{55} & \ding{51} \\
\hline
Controle de Acesso por Usuário & \ding{55} & \ding{55} & \ding{55} & \ding{55} & \ding{51} \\
\hline
\end{tabular}}
\legend{Fonte: Elaborado pelo Autor}
\end{table}